\chapter*{CONCLUSIONES, RECOMENDACIONES Y BIBLIOGRAFIA}
\addcontentsline{toc}{chapter}{CONCLUSIONES, RECOMENDACIONES Y BIBLIOGRAFIA}

% Reinicio de contadores para numeracion coherente con el capitulo 8.
\setcounter{section}{0}
\renewcommand{\thesection}{8.\arabic{section}.}
\renewcommand{\thesubsection}{8.\arabic{section}.\arabic{subsection}.}
\renewcommand{\theHsection}{anexos.\arabic{section}}
\renewcommand{\theHsubsection}{anexos.\arabic{section}.\arabic{subsection}}

\section{CONCLUSIONES}

\begin{enumerate}
  \item La vivienda unifamiliar de dos pisos de propiedad de Aquiles Taylor Ramos Yapo, ubicada en el distrito de San Miguel, provincia de San Roman y departamento de Puno, se atiende mediante una instalacion electrica residencial monofasica de 220 V. El sistema se ha sectorizado en 8 circuitos derivados independientes para garantizar selectividad y seguridad.
  \item La potencia instalada (PI) alcanza los 11.66 kW y la demanda maxima (MD) de diseno calculada para el escenario conservador (con cocina electrica) es de 10.36 kW, con una corriente demandada de 52.31 A. Esto justifica plenamente la seleccion de un alimentador general de cobre de 16 mm\(^2\), con proteccion termomagnetica de 2P-63A y proteccion diferencial de 2P-63A (30 mA).
  \item En caso de confirmarse el uso de cocina a gas en obra, la demanda maxima se reduce a 3.50 kW con una corriente demandada de 17.68 A, lo cual permitiria redimensionar de manera optima el alimentador a 10 mm\(^2\) y la proteccion general a 2P-40A.
  \item El diseno arquitectonico electrico contempla un Tablero General (TG) en el primer piso y un tablero de distribucion secundario (T2) en el segundo piso, facilitando la operacion y mantenimiento de los circuitos de cada nivel y optimizando la longitud del cableado.
  \item La seleccion de conductores (2.5 mm\(^2\) para circuitos derivados de alumbrado y tomacorrientes, 10 mm\(^2\) para cargas especiales de cocina y lavadora, e interconexion a T2) cumple con la seccion minima exigida por el Codigo Nacional de Electricidad - Utilizacion y la norma EM.010 del RNE.
  \item La incorporacion de interruptores diferenciales en todos los circuitos de tomacorrientes y fuerza, y la implementacion de un sistema de puesta a tierra (SPAT) con varilla copperweld de 5/8'' y resistencia de diseno inferior a 15 Ohmios, garantizan la seguridad e integridad de los usuarios contra contactos directos e indirectos.
\end{enumerate}

\section{RECOMENDACIONES}

\begin{enumerate}
  \item Validar fisicamente en obra las dimensiones de los ambientes y el recorrido final de las canalizaciones empotradas antes del vaciado de las losas de concreto del primer y segundo piso.
  \item Coordinar con la empresa concesionaria de energia de la region (Electro Puno) para la ubicacion definitiva del medidor de energia y la caja de toma en la fachada principal de la vivienda.
  \item Realizar una proforma comercial con las tiendas locales de distribucion (Maestro, Sodimac y Promart) para adquirir los materiales consolidados recomendados, lo que permite un ahorro significativo al comprar en empaques comerciales (rollos, tubos y cajas) segun el plan mixto optimizado de S/ 3,968.02.
  \item Efectuar la medicion de la resistencia de la puesta a tierra despues de su construccion real para validar que el valor obtenido sea menor o igual a 15 Ohmios, aplicando gel conductivo si el terreno natural de San Miguel presenta alta resistividad.
  \item Mantener rotulados de forma permanente todos los circuitos en el Tablero General (TG) y en el Tablero Secundario (T2) para facilitar la maniobra de emergencia.
\end{enumerate}

\section{BIBLIOGRAFIA}

\begin{enumerate}
  \item Ministerio de Energia y Minas. \textit{Codigo Nacional de Electricidad -- Utilizacion}. Resolucion Ministerial N. 0037-2006-MEM.
  \item Ministerio de Vivienda, Construccion y Saneamiento. \textit{Reglamento Nacional de Edificaciones (RNE)}. Norma Tecnica EM.010 Instalaciones Electricas Interiores.
  \item Instituto Nacional de Calidad (INACAL). \textit{Normas Tecnicas Peruanas (NTP) sobre Conductores Electricos e Interruptores de Proteccion}.
\end{enumerate}

\bigskip
\noindent\textbf{Nota final:} Este documento tecnico-academico ha sido elaborado con fines didacticos y de evaluacion de curso. Su implementacion real debe ser supervisada e inspeccionada en obra por un Ingeniero Mecanico Electricista o Electricista colegiado y habilitado.
